\documentclass[10pt, letterpaper]{article}

% Packages:
\usepackage[
    ignoreheadfoot, % set margins without considering header and footer
    top=1.7 cm, % seperation between body and page edge from the top
    bottom=1.8 cm, % seperation between body and page edge from the bottom
    left=2 cm, % seperation between body and page edge from the left
    right=2 cm, % seperation between body and page edge from the right
    footskip=1.0 cm, % seperation between body and footer
    % showframe % for debugging 
]{geometry} % for adjusting page geometry
\usepackage{titlesec} % for customizing section titles
\usepackage{tabularx} % for making tables with fixed width columns
\usepackage{array} % tabularx requires this
\usepackage[dvipsnames]{xcolor} % for coloring text
\definecolor{primaryColor}{RGB}{0, 79, 144} % define primary color
\usepackage{enumitem} % for customizing lists
\usepackage{fontawesome5} % for using icons
\usepackage{amsmath} % for math
\usepackage[
    pdftitle={Rohit Roy's CV},
    pdfauthor={Rohit Roy},
    pdfcreator={LaTeX with RenderCV},
    colorlinks=true,
    urlcolor=primaryColor
]{hyperref} % for links, metadata and bookmarks
\usepackage[pscoord]{eso-pic} % for floating text on the page
\usepackage{calc} % for calculating lengths
\usepackage{bookmark} % for bookmarks
\usepackage{lastpage} % for getting the total number of pages
\usepackage{changepage} % for one column entries (adjustwidth environment)
\usepackage{paracol} % for two and three column entries
\usepackage{ifthen} % for conditional statements
\usepackage{needspace} % for avoiding page brake right after the section title
\usepackage{iftex} % check if engine is pdflatex, xetex or luatex
\usepackage{multicol}
\usepackage{enumitem}
\usepackage{eso-pic}
\usepackage{atbegshi}

% Ensure that generate pdf is machine readable/ATS parsable:
\ifPDFTeX
    \input{glyphtounicode}
    \pdfgentounicode=1
    % \usepackage[T1]{fontenc} % this breaks sb2nov
    \usepackage[utf8]{inputenc}
    \usepackage{lmodern}
\fi



% Some settings:
\AtBeginEnvironment{adjustwidth}{\partopsep0pt} % remove space before adjustwidth environment
\pagestyle{empty} % no header or footer
\setcounter{secnumdepth}{0} % no section numbering
\setlength{\parindent}{0pt} % no indentation
\setlength{\topskip}{0pt} % no top skip
\setlength{\columnsep}{0cm} % set column seperation
\makeatletter
\let\ps@customFooterStyle\ps@plain % Copy the plain style to customFooterStyle
\patchcmd{\ps@customFooterStyle}{\thepage}{
    \color{gray}\textit{\small Page \thepage{} of \pageref*{LastPage}}
}{}{} % replace number by desired string
\makeatother
\pagestyle{customFooterStyle}

\titleformat{\section}{\needspace{6\baselineskip}\bfseries\large\color{primaryColor}}{}{0pt}{}[\vspace{1pt}\titlerule]

\titlespacing{\section}{
    % left space:
    -1pt
}{
    % top space:
    0.3 cm
}{
    % bottom space:
    0.2 cm
} % section title spacing

\titleformat{\subsection}{\needspace{3\baselineskip}\bfseries\normalsize\color{primaryColor}}{}{0pt}{}
\titlespacing{\subsection}{0.2 cm}{0.2 cm}{0 cm} % subsection title spacing

\renewcommand\labelitemi{$\circ$} % custom bullet points
\newenvironment{highlights}{
    \begin{itemize}[
        topsep=0.10 cm,
        parsep=0.10 cm,
        partopsep=0pt,
        itemsep=0pt,
        leftmargin=0.4 cm + 10pt
    ]
}{
    \end{itemize}
} % new environment for highlights

\newenvironment{highlightsforbulletentries}{
    \begin{itemize}[
        topsep=0.10 cm,
        parsep=0.10 cm,
        partopsep=0pt,
        itemsep=0pt,
        leftmargin=10pt
    ]
}{
    \end{itemize}
} % new environment for highlights for bullet entries


\newenvironment{onecolentry}{
    \begin{adjustwidth}{
        0.2 cm + 0.00001 cm
    }{
        0.2 cm + 0.00001 cm
    }
}{
    \end{adjustwidth}
} % new environment for one column entries

\newenvironment{subentry}{
    \begin{adjustwidth}{0.6 cm}{0.2 cm}
}{
    \end{adjustwidth}
} % new environment for one column entries under a subsection

\newenvironment{twocolentry}[2][]{
    \onecolentry
    \def\secondColumn{#2}
    \setcolumnwidth{\fill, 6.5 cm}
    \begin{paracol}{2}
}{
    \switchcolumn \raggedleft \secondColumn
    \end{paracol}
    \endonecolentry
} % new environment for two column entries

\newenvironment{header}{
    \setlength{\topsep}{0pt}\par\kern\topsep\centering\linespread{1.5}
}{
    \par\kern\topsep
} % new environment for the header

% \newcommand{\placelastupdatedtext}{% \placetextbox{<horizontal pos>}{<vertical pos>}{<stuff>}
%   \AddToShipoutPictureFG*{% Add <stuff> to current page foreground
%     \put(
%         \LenToUnit{\paperwidth-2 cm-0.2 cm+0.05cm},
%         \LenToUnit{\paperheight-1.0 cm}
%     ){\vtop{{\null}\makebox[0pt][c]{
%         \small\color{gray}\textit{Last updated in December 2025}\hspace{\widthof{Last updated in December 2025}}
%     }}}%
%   }%
% }%

\newcommand{\placelastupdatedtext}{%
  \AddToShipoutPictureFG*{%
    \put(
      \LenToUnit{\paperwidth-2.2cm},
      \LenToUnit{1.2cm}
    ){%
      \makebox[0pt][r]{%
        \small\color{gray}\textit{Last updated in October 2026}%
      }%
    }%
  }%
}


% save the original href command in a new command:
\let\hrefWithoutArrow\href

% new command for external links:
\renewcommand{\href}[2]{\hrefWithoutArrow{#1}{\ifthenelse{\equal{#2}{}}{ }{#2 }\raisebox{.15ex}{\footnotesize \faExternalLink*}}}


\begin{document}
    \newcommand{\AND}{\unskip
        \cleaders\copy\ANDbox\hskip\wd\ANDbox
        \ignorespaces
    }
    \newsavebox\ANDbox
    \sbox\ANDbox{}

    % \placelastupdatedtext
    \begin{header}
        \textbf{\fontsize{24 pt}{24 pt}\selectfont Rohit Roy}\\
        \vspace{0.1 cm}
        \fontsize{13 pt}{13 pt}\selectfont Sophia-Antipolis, France

        \vspace{0.1 cm}

        \normalsize
        % \mbox{{\color{black}\footnotesize\faMapMarker*}\hspace*{0.13cm}Your Location}%
        % \kern 0.25 cm%
        % \AND%
        % \kern 0.25 cm%
        \mbox{\hrefWithoutArrow{mailto:rohit.a.roy@inria.fr}{\color{black}{\footnotesize\faEnvelope[regular]}\hspace*{0.13cm}rohit.a.roy@inria.fr}}%
        \kern 0.25 cm%
        \AND%
        \kern 0.25 cm%
        % \mbox{\hrefWithoutArrow{tel:+91-8653837708}{\color{black}{\footnotesize\faPhone*}\hspace*{0.13cm}0541 999 99 99}}%
        % \kern 0.25 cm%
        % \AND%
        % \kern 0.25 cm%
        \mbox{\hrefWithoutArrow{https://www.rohitroy.me/}{\color{black}{\footnotesize\faLink}\hspace*{0.13cm}rohitroy.me}}%
        \kern 0.25 cm%
        \AND%
        \kern 0.25 cm%
        \mbox{\hrefWithoutArrow{https://linkedin.com/in/rohitxroy}{\color{black}{\footnotesize\faLinkedinIn}\hspace*{0.13cm}rohitxroy}}%
        \kern 0.25 cm%
        \AND%
        \kern 0.25 cm%
        \mbox{\hrefWithoutArrow{https://github.com/scarroy-02}{\color{black}{\footnotesize\faGithub}\hspace*{0.13cm}scarroy-02}}%
    \end{header}

    \vspace{0.3 cm - 0.3 cm}


    \section{About}



        
        \begin{onecolentry}
            I am a 2nd year PhD student at Centre Inria d'UniCA, working on problems in computational geometry. I previously developed a strong interest in probability theory, and my research focus has since shifted toward Topological Data Analysis and exploring topological solutions to data-driven problems.
        \end{onecolentry}

        % \vspace{0.2 cm}

        % \begin{onecolentry}
        %     The boilerplate content was inspired by \href{https://github.com/dnl-blkv/mcdowell-cv}{Gayle McDowell}.
        % \end{onecolentry}


    
    % \section{Quick Guide}

    % \begin{onecolentry}
    %     \begin{highlightsforbulletentries}


    %     \item Each section title is arbitrary and each section contains a list of entries.

    %     \item There are 7 unique entry types: \textit{BulletEntry}, \textit{TextEntry}, \textit{EducationEntry}, \textit{ExperienceEntry}, \textit{NormalEntry}, \textit{PublicationEntry}, and \textit{OneLineEntry}.

    %     \item Select a section title, pick an entry type, and start writing your section!

    %     \item \href{https://docs.rendercv.com/user_guide/}{Here}, you can find a comprehensive user guide for RenderCV.


    %     \end{highlightsforbulletentries}
    % \end{onecolentry}

    \section{Education}



        
        \begin{twocolentry}{

        \textit{Nov 2025 – Oct 2028 (expected)}}
        
            \textbf{Centre Inria d'Université Côte d'Azur}

            \textit{PhD in Computer Science and Applied Mathematics}
        \end{twocolentry}

        \vspace{0.3cm}

        \begin{twocolentry}{
            
        \textit{Aug 2023 – Apr 2025}
        
        CGPA : 9.00}
            \textbf{Chennai Mathematical Institute}

            \textit{Master of Science in Data Science}
        \end{twocolentry}

        \vspace{0.3cm}

        \begin{twocolentry}{
            
        \textit{Dec 2020 – Apr 2023}
        
        CGPA : 7.88}
            \textbf{Chennai Mathematical Institute}

            \textit{Bachelor of Science (Honours) in Mathematics and Computer Science}
        \end{twocolentry}

        % \vspace{0.3cm}

        % \begin{twocolentry}{
            
        % \textit{Apr 2018 - May 2020}
        
        % Percentage : 97.2}
        %     \textbf{DAV Model School, Durgapur}

        %     \textit{B.Sc. (Hons.) in Mathematics and Computer Science}
        % \end{twocolentry}

    % \section{Key Courses Undertaken}

    % \begin{onecolentry}
    %     \textbf{Mathematics : }
    %     \vspace{-0.5cm}
    %     \footnotesize{
    %     \begin{multicols}{3}
    %         \begin{itemize}
    %             \item Real and Complex Analysis
    %             \item Probability Theory and Stochastic Processes
    %             \item Group, Ring and Field Theory
    %             \item Linear Algebra
    %             \item Topology
    %             \item Enumerative Combinatorics
    %             \item Statistics
    %             \item Ordinary and Partial Differential Equations
    %             \item Hilbert Spaces and Quantum Probability
    %             \item Measure Theoretic Probability
    %             \item Stochastic Integration
    %         \end{itemize}
    %     \end{multicols}}
        
    % \end{onecolentry}

    % \vspace{-0.1cm}

    % \begin{onecolentry}
    %     \textbf{Computer Science and Data Science : }
    %     \vspace{-0.5cm}
    %     \footnotesize{
    %     \begin{multicols}{3}
    %         \begin{itemize}
    %             \item Design and Analysis of Algorithms
    %             \item Programming Language Concepts
    %             \item Foundations of Machine Learning
    %             \item Theory of Computation
    %             \item Topological Data Analysis
    %             \item Advanced Machine Learning
    %             \item Database Concepts
    %             \item Simulation Techniques
    %             \item Undecidability in Algebra \& Topology
    %             \item Natural Language Processing
    %             \item Applied Data Analytics
    %         \end{itemize}
    %     \end{multicols}}
        
    % \end{onecolentry}

    \section{Publications}

    \subsection{Conference Proceedings}
    \vspace{0.2cm}
    \begin{list}{}{}
      \item[\textbf{[1]}] S.~Pritam, \textbf{R.~Roy}.
      Collapse and Persistence of Directed Filtered Graphs. \href{https://epubs.siam.org/doi/10.1137/1.9781611978957.14}{DOI}\\
      \textit{SIAM Symposium on Algorithm Engineering and Experiments (ALENEX26)}
    \end{list}

    % \subsection{Journal Articles}

    % \begin{list}{}{}
    %   \item[\textbf{[n]}] Authors.
    %   Title. \href{https://doi.org/...}{DOI}\\
    %   \textit{Journal}
    % \end{list}

    \subsection{Workshop Papers}
    \vspace{0.2cm}
    \begin{list}{}{}
      \item[\textbf{[2]}] E.W.~Chambers, C.~Fillmore, S.S.~Mukherjee, \textbf{R.~Roy}, E.~Stephenson, M.~Wintraecken.
      The symmetry set as a knitting pattern for vineyards.  \\
      \textit{Computational Geometry: Young Researchers Forum (CG:YRF) 2026}

      \item[\textbf{[3]}] E.W.~Chambers, C.~Fillmore, S.S.~Mukherjee, \textbf{R.~Roy}, E.~Stephenson, M.~Wintraecken.
      Bouquet: A Visualization Tool for Symmetry Sets and Vineyards. \href{https://bouquet2d.rohitroy.me}{Tool}, \href{https://hal.science/view/index/docid/5546179}{HAL}\\
      \textit{European Workshop on Computational Geometry (EuroCG 2026)}
    \end{list}

    \subsection{Preprints}
    \vspace{0.2cm}
    \begin{list}{}{}
      \item[\textbf{[4]}] E.W.~Chambers, C.~Fillmore, S.S.~Mukherjee, \textbf{R.~Roy}, E.~Stephenson, M.~Wintraecken.
      The Singular Source of Vineyard Monodromy. \href{https://arxiv.org/abs/2607.01046}{arXiv}, \href{https://hal.science/hal-05674708v1}{HAL}

      \item[\textbf{[5]}] S.~Chowdhury, S.~Pritam, \textbf{R.~Roy}, M.C.~Sajeev.
      Filtration-Based Representation Learning for Temporal Graphs. \href{https://arxiv.org/abs/2502.10076}{arXiv}
    \end{list}

    \section{Presentations}

    \subsection{Talks}

    \vspace{0.1cm}

        \begin{subentry}
            \textit{Monodromy in Radial Persistence Vineyards and visualizing with Bouquet2D} \hfill \textit{Oct 2026}

            DataShape Seminar, Inria Sophia-Antipolis \ \href{https://www.rohitroy.me/assets/DataShape_seminar_slides.pdf}{Slides}
        \end{subentry}

        \vspace{0.2cm}

        \begin{subentry}
            \textit{Bouquet: A Visualization Tool for Symmetry Sets and Vineyards} \hfill \textit{Mar 2026}

            European Workshop on Computational Geometry (EuroCG 2026), Hagen \ \href{https://www.rohitroy.me/assets/eurocg26_slides.pdf}{Slides}
        \end{subentry}

    \subsection{Posters}

    \vspace{0.1cm}

        \begin{subentry}
            \textit{Bouquet: A Visualization Tool for Symmetry Sets and Vineyards} \hfill \textit{Sep 2026}

            AATRN Poster Session \ \href{https://www.rohitroy.me/assets/aatrn26_poster.pdf}{Poster}
        \end{subentry}

    \section{Research Experience}

        \begin{twocolentry}{
        \textit{Bordeaux, France}    
            
        \textit{May 2023 – Jul 2023}}
            \textbf{Research Intern}
            
            \textit{\href{https://www.labri.fr/en}{LaBRI}, Universite de Bordeaux}
        \end{twocolentry}

        \vspace{0.10 cm}
        \begin{onecolentry}
            \begin{highlights}
                \item Worked under \href{https://www.labri.fr/perso/marckert/}{Prof. Jean-Francois Markert} on directed animals in 2D.
                \item Investigated open combinatorial questions concerning directed animals, a concept representing percolation clusters in directed percolation models and conducted simulations of directed animals on triangular and square lattices.
                \item More details can be found \href{https://projects.lsv.ens-cachan.fr/relax/people/internships/}{here}, with the complete internship log \href{https://drive.google.com/file/d/1ipv-vD5ERt3TrnBf_6TboQ9xuqVlFqWx/view?usp=drive_link}{here}.
            \end{highlights}
        \end{onecolentry}

        \vspace{0.2cm}

        \begin{onecolentry}
            \textbf{Reading Project: Branching Processes and Convergence Concepts} \hfill \textit{May 2022 – Jul 2022}

            \textit{Under \href{https://sites.google.com/site/bhattacharya90ayan/}{Prof. Ayan Bhattacharya}, IIT Bombay}
        \end{onecolentry}

        \vspace{0.2cm}

        \begin{onecolentry}
            \textbf{Reading Project: Markov Chains and related concepts} \hfill \textit{May 2022 – Jul 2022}

            \textit{Under \href{https://www.cmi.ac.in/~rlk/}{Prof. Rajeeva Karandikar}, Chennai Mathematical Institute}
        \end{onecolentry}

    % Detailed reading project entries (replaced by the short entries above):

    % \begin{twocolentry}{
    %     \textit{Sep 2023 – Present}}
    %         \textbf{Classification of Graphs using Persistence Diagrams}
            
            
    %     \end{twocolentry}

    %     \vspace{0.10 cm}
    %     \begin{onecolentry}
    %         \begin{highlights}
    %             \item Part of a grant funded by Fujitsu Research, working under \href{https://siddharthpritam.github.io/}{Prof. Siddharth Pritam} at Chennai Mathematical Institute.
    %             \item The goal of this project is to classify temporal graphs using methods from TDA. We extract topological features from these graphs and then these are fed into kernel-based methods for classification.
    %             \item The arXiv preprint for this project is \href{https://arxiv.org/abs/2502.10076}{here.}
    %         \end{highlights}
    %     \end{onecolentry}

    % \vspace{0.2cm}

    % \begin{twocolentry}{
    % \textit{May 2022 – Jul 2022}}
    % \textbf{Branching Processes and Convergence Concepts}
            
            
    % \end{twocolentry}

    % \vspace{0.10 cm}
    % \begin{onecolentry}
    % \begin{highlights}
    % \item This was a reading project under \href{https://sites.google.com/site/bhattacharya90ayan/}{Prof. Ayan Bhattacharya} from IIT Bombay.
    % \item Learnt about Galton-Watson process, genealogical trees and the probability of
    % extinction of the same process.
    % \item Discussed about various limit theorems in probability and had a brief overview of branching random walks.
    % \end{highlights}
    % \end{onecolentry}

    % \vspace{0.2cm}

    % \begin{twocolentry}{
    % \textit{May 2022 – Jul 2022}}
    % \textbf{Markov Chains and related concepts}
            
            
    % \end{twocolentry}

    % \vspace{0.10 cm}
    % \begin{onecolentry}
    % \begin{highlights}
    % \item This was a reading project under \href{https://www.cmi.ac.in/~rlk/}{Prof. Rajeeva Karandikar} from Chennai Mathematical Institute.
    % \item Guided reading of the book ``An Introduction to Probability Theory and its Applications" by William Feller.
    % \item Learnt about Markov Chains and Stochastic Processes.
    % \end{highlights}
    % \end{onecolentry}

        % \vspace{0.2cm}
    % \section{Publications}



        
    %     \begin{samepage}
    %         \begin{twocolentry}{
    %             Jan 2004
    %         }
    %             \textbf{3D Finite Element Analysis of No-Insulation Coils}

    %             \vspace{0.10 cm}

    %             \mbox{Frodo Baggins}, \mbox{\textbf{\textit{John Doe}}}, \mbox{Samwise Gamgee}
    %         \end{twocolentry}


    %         \vspace{0.10 cm}

    %         \begin{onecolentry}
    %     \href{https://doi.org/10.1109/TASC.2023.3340648}{10.1109/TASC.2023.3340648}
    %         \end{onecolentry}
    %     \end{samepage}

    

    \section{Teaching Experience}

    \begin{twocolentry}{
        \textit{Chennai, India}}
            \textbf{Chennai Mathematical Institute}

            \textit{Teaching Assistant for the following courses}
    \end{twocolentry}

    \begin{onecolentry}
        \begin{highlights}
            \item Data Mining and Machine Learning \hfill \textit{Jan 2025 – Apr 2025}\\
            {\small Instructor: Prof. Madhavan Mukund \quad·\quad For: M.Sc. Data Science 1st Year}
            \item Distributed Computing and Big Data \hfill \textit{Jan 2025 – Apr 2025}\\
            {\small Instructor: Prof. Venkatesh Vinayakrao \quad·\quad For: M.Sc. Data Science 1st Year}
            \item Introduction to Martingales \hfill \textit{Jan 2025 – Apr 2025}\\
            {\small Instructor: Prof. Rajeeva Karandikar \quad·\quad For: B.Sc. 3rd Year and M.Sc.}
            \item Natural Language Processing \hfill \textit{Aug 2024 – Dec 2024}\\
            {\small Instructor: Prof. Ramaseshan Ramachandran \quad·\quad For: M.Sc. Data Sci. \& M.Sc. Comp. Sci.}
            \item Probability Theory \hfill \textit{Jan 2024 – Apr 2024}\\
            {\small Instructor: Prof. Rajeeva Karandikar \quad·\quad For: B.Sc. Mathematics 1st Year}
            \item Mathematical Methods in Analysis \hfill \textit{Aug 2022 – Nov 2022}\\
            {\small Instructor: Prof. Rajeeva Karandikar \quad·\quad For: M.Sc. Data Science 1st Year}
        \end{highlights}
    \end{onecolentry}

    \section{Academic and Volunteer Service}

    \begin{onecolentry}
        \begin{highlights}
            \item \textbf{Artifact Evaluation Committee member}, SIAM Symposium on Algorithm Engineering and Experiments (ALENEX 2027).
        \end{highlights}
    \end{onecolentry}

    \section{Awards}

    \begin{onecolentry}
        \begin{highlights}
            \item \textbf{Olympiad Mathematics :} Qualified for the Indian National Mathematical Olympiad (INMO) in 2018.
            \item \textbf{KVPY-SA 2018 Fellow :} Recipient of the KVPY scholarship given yearly. 
            \item \textbf{Mathematics Talent Reward Programme (MTRP)} held in ISI Kolkata : Qualified in 2017 and 2019, and received 3rd prize in 2017. 
        \end{highlights}
    \end{onecolentry}


    \section{Industry Experience}

        \begin{twocolentry}{
        \textit{Bengaluru, India}    
            
        \textit{Jun 2025 – Oct 2025}}
            \textbf{Machine Learning Engineer}
            
            \textit{Hewlett-Packard, Printing and PC Systems (PPS)}
        \end{twocolentry}

        \vspace{0.10 cm}
        \begin{onecolentry}
            \begin{highlights}
                \item Contributed to the development of an enterprise-scale question–answering system for printing solutions at HP.
                \item Extended a document-based language model pipeline to support multimodal inputs, incorporating information from images and tabular data.
                \item Implemented and evaluated models and workflows on the AWS SageMaker platform.
                \item Investigated and applied prompt engineering techniques to improve model interpretability and response quality.
                \item Designed mechanisms for structured metadata retrieval and document access from SharePoint to support efficient information retrieval.
            \end{highlights}
        \end{onecolentry}

        \vspace{0.2 cm}

        \begin{twocolentry}{
        \textit{Chennai, India}    
            
        \textit{May 2024 – Aug 2024}}
            \textbf{Data Analysis Intern (Industry Internship)}
            
            \textit{\href{https://nptel.ac.in/}{NPTEL}, IIT Madras}
        \end{twocolentry}

        \vspace{0.10 cm}
        \begin{onecolentry}
            \begin{highlights}
                \item Classified colleges on their performance in NPTEL courses using clustering and came up with metrics for comparison.
                \item Created System Dynamic (SD) models using Vensim to mimic NPTEL course enrollment and registration trends across semesters with real world factors in play.
            \end{highlights}
        \end{onecolentry}
    
    \section{Languages}

       \begin{onecolentry}
           \textbf{Bengali :} Native
           \hspace{1cm}
           \textbf{English :} C2
           \hspace{1cm}
           \textbf{Hindi :} C2
           \hspace{1cm}
           \textbf{French :} A1
           \hspace{1cm}
           \textbf{German :} A1
       \end{onecolentry}

\placelastupdatedtext
    

\end{document}